% SOURCE_AUTHORITY: sga5_fr_workpass.tex, lines 14011--14038 (LF-normalized unit).
% SOURCE_SHA256: 791F4EFFC5E02832D5D77ED03518C8156D6F07E4C8238B03545DB93D883FBB28
\subsubsection*{b)} Seja $x'\in X_x^{\prime f'}$ e seja $G_{x'}$ o estabilizador do ponto $x'$. Evidentemente, $G_{x'}$ é estável por $\varphi$, obtendo-se, portanto, um endomorfismo
\[
\varphi_{x'} = \varphi|G_{x'} : G_{x'} \longrightarrow G_{x'} \quad .
\]

Denotemos por $X'/G_{x'}$ o esquema quociente obtido a partir de $X'$ pela relação de equivalência definida por $G_{x'}$, por $f'/G_{x'}$ o endomorfismo de $X'/G_{x'}$ obtido a partir do endomorfismo $f'$ de $X'$ mediante a mesma relação de equivalência e por $x'$ a classe do ponto $x'$ módulo essa relação de equivalência.% SOURCE_ANOMALY_PRESERVED: source prints capital F'/G although it denotes the descended endomorphism f'/G.

A projeção canônica
\[
p : X' \longrightarrow X'/G = X
\]
fatora-se da seguinte maneira:
\[
p : X' \xrightarrow{\;p_1\;} X'/G_{x'} \xrightarrow{\;p_2\;} X = X'/G \quad .
\]

Tomando a imagem inversa por $p_2$ do homomorfismo $u$, obtém-se um homomorfismo
\[
v = p_2(u) : (f'/G_{x'})^*(p_2^*(F)) \longrightarrow p_2^*(F) \quad .
\]

É fácil ver que:

\begin{equation}\tag{6.4.3}
\alpha_{x'}(f'/G_{x'}, v, f', G_{x'}, \varphi_{x'}) = \alpha_x(f,u, f', G, \varphi) \quad (x = p(x')) \quad .
\end{equation}

Pode-se, portanto, sempre reduzir ao caso «totalmente ramificado», isto é, ao caso em que a fibra do ponto $x$ se reduz a um ponto.
